\exercisesection{复合积}

\begin{enumerate}
\def\labelenumi{\arabic{enumi})}
\tightlist
\item
  设 \(k\) 是一个交换环，A 是一个带单位的结合 \(k\) -代数，M 是一个 (A, A)-双模。A 被 M 的一个扩张是如下数据：一个带单位的结合 \(k\) -代数 B 以及一个 \(k\) -模的正合列
\end{enumerate}

\[
(\mathcal {E}): 0 \rightarrow \mathrm{M} \stackrel {{i}} {{\rightarrow}} \mathrm{B} \stackrel {{p}} {{\rightarrow}} \mathrm{A} \rightarrow 0,
\]

其中 p 是一个代数同态，并且

\[
i (p (b) m) = b i (m) \quad \text { and } \quad i (m p (b)) = i (m) b \quad \text { for } \quad b \in \mathbf {B}, \quad m \in \mathbf {M}.
\]

两个扩张 \(0 \to M \xrightarrow{i} B \xrightarrow{p} A \to 0\) \(0 \to M \xrightarrow{i'} B' \xrightarrow{p'} A \to 0\) 如果存在一个代数同态，则称它们是等价的 \(u: B \to B'\) 使得 \(p' \circ u = p\) \(u \circ i = i'\) ；同态 \(u\) 则它是双射的。

\begin{enumerate}
\def\labelenumi{\alph{enumi})}
\item
  我们假设一个 \(k\) -模的正合列 \(0 \to \mathbf{M} \xrightarrow{i} \mathbf{B} \xrightarrow{p} \mathbf{A} \to 0\) 是分裂的；令 \(s: \mathbf{A} \to \mathbf{B}\) 为一个 \(k\) -线性映射，它是 \(p\) 的一个截面。证明存在一个元素 \(f \in \mathbf{C}^2(\mathbf{A}, \mathbf{M})\) （X，第 195 页，习题 26，b））使得 \(s(a) s(a') - s(aa') = i(f(a, a'))\) （当 \(a, a' \in \mathbf{A}\) 时）。证明 \(f\) 是一个 2-上循环，且其类在 \(\mathbf{H}^2(\mathbf{A}, \mathbf{M})\) 中与 \(s\) 的选择无关；我们将其记为 \(c(\mathcal{E})\) 。
\item
  设 \(\mathrm{Ex}_{k}(\mathbf{A},\mathbf{M})\) 为 A 被 M 的扩张的等价类组成的集合，这些扩张作为 k-模的扩张是平凡的。证明 c 定义了从 \(\mathrm{Ex}_{k}(\mathbf{A},\mathbf{M})\) 到 \(\mathrm{H}^{2}(\mathbf{A},\mathbf{M})\) 上的一个双射。
\item
  设 \(\mathcal{E}:0\to\mathbf{M}\stackrel{i}{\to}\mathbf{B}\stackrel{p}{\to}\mathbf{A}\to0\) 与 \(\mathcal{E}':0\to\mathbf{M}\stackrel{i'}{\to}\mathbf{B}\stackrel{p'}{\to}\mathbf{A}\to0\) 为 A 被 M 的两个扩张。我们用 C 表示 \(\mathbf{B}\times\mathbf{B}'\) 的由元素 \((b,b')\) 组成的子代数，这些元素满足 \(p(b)=p'(b')\) ；用 c 表示 C 的由元素 \((i(m),-i'(m))\) （\(m\in\mathbf{M}\)）组成的理想，用 \(\mathbf{B}''\) 表示代数 C/c，用 \(\pi:\mathbf{C}\to\mathbf{B}''\) 表示商同态。设 \(i'':\mathbf{M}\to\mathbf{B}''\) 与 \(p'':\mathbf{B}''\to\mathbf{A}\) 为满足 \(i''(m)=\pi((m,0))\) （\(m\in\mathbf{M}\)）与 \(p''\pi((b,b'))=p(b)\) （\(b\in\mathbf{B},b'\in\mathbf{B}'\)）的同态。证明序列 \(0\to\mathbf{M}\to\mathbf{B}''\to\mathbf{A}\to0\) 是 A 被 M 的一个扩张，称为 ( \(\mathcal{E}\) ) 与 ( \(\mathcal{E}'\) ) 的和扩张，并且这在 A 被 M 的扩张的等价类集合上定义了一个交换群法则。特别地，集合 \(\mathrm{Ex}_{k}(\mathbf{A},\mathbf{M})\) 由此被赋予一个群结构，且映射 \(c\) 是一个群同构。
\end{enumerate}

\begin{enumerate}
\def\labelenumi{\arabic{enumi})}
\setcounter{enumi}{1}
\item
  假设 \(k\) -代数 A 具有一个增广 \(\pi:\mathbf{A}\to k\) （X，第 196 页，习题 27），并且
\item
  假设 \(k\) -代数 \(\mathbf{A}\) 允许一个增广 \(\pi : \mathbf{A} \to k\) （X，第 196 页，习题 27），而且此外， \(\mathbf{A}\) 是一个投射 \(k\) -模。我们用 I 表示 \(\pi\) 的核。设 \(\mathbf{M}\) 为一个左 \(\mathbf{A}\) -模。
\end{enumerate}

\begin{enumerate}
\def\labelenumi{\alph{enumi})}
\item
  设 \((\mathcal{E}):0\to M_{(\pi)}\stackrel{i}{\to}\mathbf{B}\stackrel{p}{\to}\mathbf{A}\to 0\) 为 \(\mathbf{A}\) 被 \((\mathbf{A},\mathbf{A})\) -双模 \(M_{(\pi)}\) 的一个扩张（X，第 196 页，习题 27）；我们用 \(\overline{\mathbf{B}}\) 表示 \(\mathbf{B}\) 的由元素 \(b\in \mathbf{B}\) 组成的理想，这些元素满足 \(p(b)\in I\) 。我们有 \(i(m)b = 0\) ，对 \(m\in \mathbf{M}\) 与 \(b\in \overline{\mathbf{B}}\) ；于是 \(\overline{\mathbf{B}}\) 具有一个左 \(\mathbf{A}\) -模结构，使得 \(p(\beta)b = \beta b\) 对 \(\beta \in \mathbf{B}\) , \(b\in \overline{\mathbf{B}}\) 成立。证明存在一个 \(\mathbf{A}\) -模的正合列 \(0\to \mathrm{M}\to \overline{\mathrm{B}}\to I\to 0\) ；我们用 \(\theta (\mathcal{E})\) 表示该正合列在 \(\operatorname {Ext}_{\mathbf{A}}^{1}(I,\mathbf{M})\) 中的类。
\item
  证明 \(\theta\) 定义了一个从 \(\mathrm{Ex}_{k}(\mathbf{A},\mathbf{M}_{(\pi)})\) 到 \(\mathrm{Ext}_{\mathbf{A}}^{1}(I,\mathbf{M})\) 的映射。证明图表的交换性：
\end{enumerate}

\[
\begin{array}{c c c} \operatorname{Ex} _ {k} (\mathbf {A}, \mathbf {M} _ {(\pi)}) & \stackrel {{\theta}} {{\longrightarrow}} & \operatorname{Ext} _ {\mathbf {A}} ^ {1} (I, \mathbf {M}) \\ \downarrow^ {c} & & \downarrow^ {- \delta} \\ \mathbf {H} ^ {2} (\mathbf {A}, \mathbf {M} _ {(\pi)}) & \stackrel {{\varphi}} {{\longrightarrow}} & \operatorname{Ext} _ {\mathbf {A}} ^ {2} (k, \mathbf {M}) \end{array}
\]

其中 δ 是由正合列 \(0 \rightarrow I \rightarrow A \xrightarrow{\pi} k \rightarrow 0\) 推出的连接同态，而 φ 是借助标准分解进行计算而得到的同构（X，第 196 页，习题 27，b））。由此推出 θ 是群的同构。

\begin{enumerate}
\def\labelenumi{\arabic{enumi})}
\setcounter{enumi}{2}
\tightlist
\item
  设 G 为一个群，A 为代数 \(\mathbf{Z}^{(\mathrm{G})}, j: \mathrm{G} \to \mathrm{A}\) ，j 为典范单射。我们以由 \(\pi: A \to Z\) （\(\pi(e_g) = 1\) 对所有 \(g \in G\) ）定义的增广来赋予 A。设 M 为一个左 A-模。
\end{enumerate}

\begin{enumerate}
\def\labelenumi{\alph{enumi})}
\item
  设 \(0 \to M_{(\pi)} \xrightarrow{i} B \xrightarrow{p} A \to 0\) 为 \(A\) 被 \((A, A)\) -双模 \(M_{(\pi)}\) 的一个扩张（X，第 196 页，习题 27）。我们用 \(\Gamma\) 表示 \(B\) 中由元素 \(b\) （属于 \(B\) 且满足 \(p(b) \in j(G)\) ）组成的子集。证明 \(\Gamma\) 赋予由 \(B\) 的乘法诱导的乘法后，是一个群，即 \(G\) 被 \(G\) -模 \(M\) 的一个扩张。
\item
前面的构造定义了一个从 \(\mathrm{Ex}_{\mathbf{Z}}(\mathbf{A},\mathbf{M}_{(\pi)})\) 到 \(\mathrm{Ex}(\mathbf{G},\mathbf{M})\) 的映射（VIII，§ 13）。证明图表
\end{enumerate}

\[
\begin{array}{c} \operatorname{Ex} _ {\mathbf {Z}} (\mathbf {A}, \mathbf {M} _ {(\pi)}) \to \operatorname{Ex} (\mathbf {G}, \mathbf {M}) \\ \downarrow^ {c} \qquad \qquad \qquad \qquad \qquad \qquad \qquad \qquad \qquad \qquad \downarrow \\ \mathbf {H} ^ {2} (\mathbf {A}, \mathbf {M} _ {(\pi)}) \xrightarrow {\varphi^ {2}} \mathbf {H} ^ {2} (\mathbf {G}, \mathbf {M}) \end{array}
\]

是交换的。

* 4) 设 g 为一个李 k-代数，M 为一个 g-模。g 被 M 的扩张是 k-模的一个正合列 \(0 \rightarrow M \xrightarrow{i} h \xrightarrow{p} g \rightarrow 0\) ，其中 h 是李 k-代数，p 是李代数同态，且 \(i(p(H)m) = [H, i(m)]\) 对 \(H \in h, m \in M\) 成立。两个扩张

\[
0 \rightarrow \mathrm{M} \stackrel {i} {\rightarrow} \mathfrak {h} \stackrel {p} {\rightarrow} \mathfrak {g} \rightarrow 0, 0 \rightarrow \mathrm{M} \stackrel {i ^ {\prime}} {\rightarrow} \mathfrak {h} ^ {\prime} \stackrel {p ^ {\prime}} {\rightarrow} \mathfrak {g} \rightarrow 0
\]

被称为等价的，如果存在一个同态 \(u: \mathfrak{h} \to \mathfrak{h}'\) 使得 \(p' \circ u = p\) , \(u \circ i = i'\) 。我们用 Liex (g, M) 表示 g 被 M 扩张的等价类集合。我们假设 g 是一个自由 k-模。

\begin{enumerate}
\def\labelenumi{\alph{enumi})}
\tightlist
\item
  设 \((\mathcal{E}):0\to \mathrm{M}\stackrel {i}{\to}\mathfrak{h}\stackrel {p}{\to}\mathfrak{g}\to 0\) 为 \(\mathfrak{g}\) 被 \(\mathbf{M}\) 的一个扩张；我们选择一个 \(k\) -线性截面 \(s:\mathfrak{g}\to \mathfrak{h}\) （相对于 \(p\) ）。我们用下式定义一个元素 \(f\in C^2 (\mathfrak{g},\mathbf{M})\) （X，第 196 页，习题 28，b））：
\end{enumerate}

\[
i (f (x, y)) = [ s (x), s (y) ] - s ([ x , y ]) \quad \text { for } \quad x  ,   y \in \mathfrak {g}  .
\]

证明 \(f\) 是一个 2-上循环，且它在 \(\mathbf{H}^{2}(\mathfrak{g},\mathbf{M})\) 中的类不依赖于 \(s\) 的选择；我们将其记作 \(\theta(\mathcal{E})\) 。b) 证明 \(\theta\) 定义了从 Liex \((\mathfrak{g},\mathbf{M})\) 到 \(\mathbf{H}^{2}(\mathfrak{g},\mathbf{M})\) 上的一个双射。直接描述 Liex \((\mathfrak{g},\mathbf{M})\) 上经由 \(\gamma\) 转移而得到的群结构。

\begin{enumerate}
\def\labelenumi{\alph{enumi})}
\setcounter{enumi}{2}
\tightlist
\item
  设 U 为 g 的包络代数， \(M_{(\pi)}\) 为借助增广 \(\pi: U \to k\) 与 M 相关联的 (U, U)-双模（X，第 196 页，习题 27）。设 \(0 \to M_{(\pi)} \to B \xrightarrow{p} U \to 0\) 为 U 被 \(M_{(\pi)}\) 的一个扩张（习题 1）；我们用 b 表示 B 中满足 \(p(b) \in j(g)\) 的元素 b 的集合。
\end{enumerate}

的元素b的集合。证明b是一个李k-代数，是g通过M的一个扩张。此构造定义了一个映射 \(\lambda: \mathrm{Ex}_{k}(\mathrm{U}, \mathrm{M}_{(\pi)}) \to \mathrm{Liex}(\mathfrak{g}, \mathrm{M})\) 。证明该图表

\[
\begin{array}{c c c} \operatorname{Ex} _ {k} (\mathrm{U}, \mathrm{M} _ {(\pi)}) & \xrightarrow {\lambda} & \text { Liex(g,M) } \\ \downarrow^ {c} & & \downarrow^ {\gamma} \\ \mathrm{H} ^ {2} (\mathrm{U}, \mathrm{M} _ {(\pi)}) & \xrightarrow {\varphi^ {2}} & \mathrm{H} ^ {2} (\mathrm{g,M}) \end{array}
\]

（其中 \(\varphi^2\) 是 X 第 196 页习题 27 所定义的同构）是交换的。由此推出 \(\lambda\) 是一个同构。*

¶ 5) 设G、F为两个群。

\begin{enumerate}
\def\labelenumi{\alph{enumi})}
\item
  对于每个扩张 \(\mathcal{E}: F \to E \to G\) （即 \(G\) 被 \(F\) 的扩张）(I，第 62 页)， \(E\) 借助内自同构对 \(F\) 的作用定义了一个同态 \(E \to \text{Aut}(F)\) ，从而通过取商得到一个群同态 \(\theta: G \to \text{Aut}(F)/\text{Int}(F)\) ，称为与扩张 \(\mathcal{E}\) 相伴的同态。 \(G\) 被 \(F\) 的两个同构的扩张具有相同的相伴同态。
\item
  反之，我们考虑一个群同态 \(\theta: G \to \text{Aut} (F)/\text{Int} (F)\) 。设
\end{enumerate}

\[
p: \operatorname{Aut} (\mathbf {F}) \rightarrow \operatorname{Aut} (\mathbf {F}) / \operatorname{Int} (\mathbf {F})
\]

商映射。我们选择一个从 G 到 Aut(F) 的映射 \(\sigma\) ，使得 \(p \circ \sigma = \theta\) ，以及一个映射 \(f\) ，从 G \(\times\) G 到 F，使得 \(\sigma(x)\, \sigma(y)\, (\sigma(xy))^{-1} = \text{Int}(f(x,y))\) 对 \(x,y \in G\) 成立；令 \(c(x,y,z) = \sigma(x)\, (f(y,z)).f(x,yz).(f(xy,z))^{-1}.(f(x,y))^{-1}\) 。证明 \(c(x,y,z)\) 属于 F 的中心 Z。

\begin{enumerate}
\def\labelenumi{\alph{enumi})}
\setcounter{enumi}{2}
\item
  我们通过令 \(g . z = \sigma(g) (z)\) （\(g \in G, z \in Z\) ）而将 Z 视为 G-模；此定义与 \(\sigma\) 的选择无关。映射 c 定义了 \(C^{3}(G, Z)\) 的一个元素；证明 c 是一个 3-上循环，且它在 \(H^{3}(G, Z)\) 中的类与 \(\sigma\) 以及 f 的选择无关。我们将其记作 \(\omega(G, F, \theta)\) 。d) 证明 \(\omega(G, F, \theta)\) 的消失是存在 G 被 F 的扩张且其相伴同态为 \(\theta\) 的充分必要条件。若 \(\omega(G, F, \theta) = 0\) ，证明可以选择 f 使得上循环 c 为零，然后在 \(G \times F\) 上定义一个群法则。
\item
  假设 \(\omega(G, F, \theta) = 0\) 。证明 \(G\) 被 \(F\) 的、相伴同态为 \(\theta\) 的扩张的同构类集合是群 H \(^{2}\) (G, Z) 作用下的一个主齐性集合。
\item
  设 H 为一个群，M 为一个 \(\mathbf{Z}^{(H)}\) -模， \(x\) 为 \(\mathrm{H}^3 (\mathrm{H},\mathrm{M})\) 的一个元素。证明存在一个以 M 为中心的群 E 以及一个同态 \(\varphi :\mathrm{H}\to \mathrm{Aut}(\mathrm{E}) / \mathrm{Int}(\mathrm{E})\) ，它在 M 上诱导出给定的 H-模结构，使得 \(\omega (\mathrm{H},\mathrm{E},\varphi) = x\) （取 \(\mathbf{E} = \mathbf{M}\times \mathbf{L}\) ，其中 L 是以 \(\mathrm{H}\times \mathrm{H}\) 为生成集的自由群）.
\end{enumerate}

\begin{enumerate}
\def\labelenumi{\arabic{enumi})}
\setcounter{enumi}{5}
\tightlist
\item
  设 M 为一个右 A-模，N 为一个左 A-模。对 \(n \geqslant 0\) ，我们用 \(\mathcal{T}_{n}(M, N)\) 表示三元组 \((P, p, q)\) 的集合，其中 P 是由有限生成自由 A-模构成的复形，使得 \(P_{r} = 0\) 对 \(r < 0\) 与 \(r > n\) 成立，且 \(p: P \to N\) 与 \(q: P^{*} \to M(n)\) 是 A-复形的态射（我们用 \(P^{*}\) 表示 Homgr 复形 \(_{A}\) \((P, A)\) ）。两个元素 \((P, p, q)\) 与 \((P', p', q')\) （属于 \(\mathcal{T}_{n}(M, N)\) ）称为相连的，如果存在态射 \(u: P \to P'\) 使得 \(p = p' \circ u\) 与 \(q' = q \circ 'u\) ；我们用 \(T_{n}(M, N)\) 表示 \(\mathcal{T}_{n}(M, N)\) 对于使两个相关元素等价的最细等价关系所作的商（参见 II，第 52 页，习题 9）。a) 证明 \(T_{n}(M, N)\) 与 Tor \(_{n}^{A}\) \((M, N)\) 等同。
\end{enumerate}

\begin{enumerate}
\def\labelenumi{\alph{enumi})}
\setcounter{enumi}{1}
\tightlist
\item
  设 \((\mathcal{E}):0\to \mathrm{M}^{\prime}\stackrel {i}{\to}\mathrm{M}\stackrel {\pi}{\to}\mathrm{M}^{\prime \prime}\to 0\) 为右 A-模的一个正合列。描述连接同态 \(\partial :\mathbf{T}_n(\mathbf{M}'',\mathbf{N})\to \mathbf{T}_{n - 1}(\mathbf{M}',\mathbf{N})\) ，它由 \(\partial (\mathcal{E},\mathbf{N})\) 利用之前的等同关系而得到。
\end{enumerate}

\((\mathrm{If}(\mathbf{P},p,q)\in\mathcal{T}_{n}(\mathbf{M}^{\prime\prime},\mathbf{N}),\text{证明我们有}\partial(\mathbf{P},p,q)=(\overline{\mathbf{P}},\overline{p},\overline{q}),\text{其中}\overline{\mathbf{P}}\text{ 是将 } \mathbf{P}_{n}\text{ 替换为 0 后得到的 } \mathbf{P},\overline{p}\text{ 是由 } p\text{ 导出的态射，而 }(\overline{q})^{n-1}: \mathbf{P}_{n-1}^{*}\to\mathbf{M}^{\prime}\text{ 等于 }u^{n-1},u\text{ 是从 } \mathbf{P}^{*}\text{ 到复形 (Con (i)) (n) 中的态射，使得 }\pi\circ u^{n}=q^{n}.)\)

¶ 7) 设 A' 是一个（结合且含幺的）k-代数，且 B = A ⊗k A'。

\begin{enumerate}
\def\labelenumi{\alph{enumi})}
\tightlist
\item
  设 M 为左 A-模，Q 为右 A-模，M' 为左 A'-模，Q' 为右 A'-模。设 (P, p)（相应地，(P', p')，(R, r)）为 A-模 M（相应地，A'-模 M'，B-模 M ⊗\_k M'）的一个投射分解。证明存在一个 B-复形的态射
\end{enumerate}

\[
u: \mathbf {P} \otimes_ {k} \mathbf {P} ^ {\prime} \rightarrow \mathbf {R},
\]

在同伦意义下唯一，使得 \(r \circ u = p \otimes p'\) 。利用 \(u\) ，定义一个零次的分次 \(k\) -同态

\[
\top : \operatorname{Tor} ^ {\mathbf {A}} (\mathrm{Q}, \mathrm{M}) \otimes_ {k} \operatorname{Tor} ^ {\mathbf {A} ^ {\prime}} (\mathrm{Q} ^ {\prime}, \mathrm{M} ^ {\prime}) \rightarrow \operatorname{Tor} ^ {\mathbf {B}} (\mathrm{Q} \otimes_ {k} \mathrm{Q} ^ {\prime}, \mathrm{M} \otimes_ {k} \mathrm{M} ^ {\prime}).
\]

证明 \(\top\) 不依赖于分解 P、P'、R 以及态射 \(u\) 的选择。如果 \(a \in \operatorname{Tor}^{\mathbf{A}}(\mathbf{Q}, \mathbf{M})\) 且 \(a' \in \operatorname{Tor}^{\mathbf{A}'}(\mathbf{Q}', \mathbf{M}')\) ，我们置 \(\top(a \otimes a') = a \top a'\) 。

\begin{enumerate}
\def\labelenumi{\alph{enumi})}
\setcounter{enumi}{1}
\tightlist
\item
  我们进一步假设 \(k\) -模 \(\mathbf{A}\) 与 \(\mathbf{A}'\) 是投射的，并且有 \(\operatorname{Tor}_n^k (\mathbf{M},\mathbf{M}') = 0\) （其中 \(n > 0\) ）。证明此时 \(\mathbf{P} \otimes_k \mathbf{P}'\) 是 \(\mathbf{B}\) -模 \(\mathbf{M} \otimes_k \mathbf{M}'\) 的一个投射分解；由此推出，对每个 \(\mathbf{A}\) -模 \(\mathbf{N}\) 与每个 \(\mathbf{A}'\) -模 \(\mathbf{N}'\) ，存在一个零次的分次 \(k\) -同态
\end{enumerate}

\[
\vee : \operatorname{Ext} _ {\mathbf {A}} (\mathbf {M}, \mathbf {N}) \otimes_ {k} \operatorname{Ext} _ {\mathbf {A} ^ {\prime}} (\mathbf {M} ^ {\prime}, \mathbf {N} ^ {\prime}) \rightarrow \operatorname{Ext} _ {\mathbf {B}} (\mathbf {M} \otimes_ {k} \mathbf {M} ^ {\prime}, \mathbf {N} \otimes_ {k} \mathbf {N} ^ {\prime}).
\]

证明 \(\vee\) 与分解 \(\mathbf{P}\) 和 \(\mathbf{P}'\) 的选择无关。如果 \(b \in \operatorname{Ext}_{\mathbf{A}}(\mathbf{M}, \mathbf{N})\) 且 \(b' \in \operatorname{Ext}_{\mathbf{A}'}(\mathbf{M}', \mathbf{N}')\) ，我们记 \(\vee (b \otimes b') = b \vee b'\) 。

\begin{enumerate}
\def\labelenumi{\alph{enumi})}
\setcounter{enumi}{2}
\tightlist
\item
  设 A'' 是第三个 k-代数（结合的且有单位元），M'' 和 N'' 是左 A''-模，Q'' 是右 A''-模。我们将 k-代数 A ⊗\_k (A' ⊗\_k A'') 与 (A ⊗\_k A') ⊗\_k A'' 等同起来，同样将 k-模 M ⊗\_k (M' ⊗\_k M'') 与 (M ⊗\_k M') ⊗\_k M''，N ⊗\_k (N' ⊗\_k N'') 与 (N ⊗\_k N') ⊗\_k N''，Q ⊗\_k (Q' ⊗\_k Q'') 与 (Q ⊗\_k Q') ⊗\_k Q'' 等同起来。证明这些公式
\end{enumerate}

\[
\begin{array}{l l} a \top (a ^ {\prime} \top a ^ {\prime \prime}) = (a \top a ^ {\prime}) \top a ^ {\prime \prime} & \text {for} \quad a \in \operatorname{Tor} ^ {\mathbf {A}} (\mathrm{Q}, \mathrm{M}), \quad a ^ {\prime} \in \operatorname{Tor} ^ {\mathbf {A} ^ {\prime}} (\mathrm{Q} ^ {\prime}, \mathrm{M} ^ {\prime}), \\ & \qquad \qquad \qquad \qquad \qquad \qquad \qquad \qquad \qquad \qquad \qquad \qquad \qquad \qquad a ^ {\prime \prime} \in \operatorname{Tor} ^ {\mathbf {A} ^ {\prime \prime}} (\mathrm{Q} ^ {\prime \prime}, \mathrm{M} ^ {\prime \prime}) \end{array}
\]

\[
\begin{array}{r l} b \vee (b ^ {\prime} \vee b ^ {\prime \prime}) = (b \vee b ^ {\prime}) \vee b ^ {\prime \prime} & \text {for} \quad b \in \operatorname{Ext} _ {\mathbf {A}} (\mathbf {M}, \mathbf {N}), \quad b ^ {\prime} \in \operatorname{Ext} _ {\mathbf {A}} (\mathbf {M} ^ {\prime}, \mathbf {N} ^ {\prime}), \\ & \qquad \qquad \qquad \qquad \qquad \qquad b ^ {\prime \prime} \in \operatorname{Ext} _ {\mathbf {A} ^ {\prime \prime}} (\mathbf {M} ^ {\prime \prime}, \mathbf {N} ^ {\prime \prime}). \end{array}
\]

\begin{enumerate}
\def\labelenumi{\alph{enumi})}
\setcounter{enumi}{3}
\tightlist
\item
  设 \(\sigma : \mathrm{Tor}^{\mathbf{A} \otimes_{\mathbf{k}} \mathbf{A}'}(\mathbf{Q} \otimes_{\mathbf{k}} \mathbf{Q}', \mathbf{M} \otimes_{\mathbf{k}} \mathbf{M}') \to \mathrm{Tor}^{\mathbf{A}' \otimes_{\mathbf{k}} \mathbf{A}}(\mathbf{Q}' \otimes_{\mathbf{k}} \mathbf{Q}, \mathbf{M}' \otimes_{\mathbf{k}} \mathbf{M})\) 与
\end{enumerate}

\[
\tau : \operatorname{Ext} _ {\mathbf {A} \otimes_ {k} \mathbf {A} ^ {\prime}} (\mathrm{M} \otimes_ {k} \mathrm{M} ^ {\prime}, \mathrm{N} \otimes_ {k} \mathrm{N} ^ {\prime}) \rightarrow \operatorname{Ext} _ {\mathbf {A} ^ {\prime} \otimes_ {k} \mathbf {A}} (\mathrm{M} ^ {\prime} \otimes_ {k} \mathrm{M}, \mathrm{N} ^ {\prime} \otimes_ {k} \mathrm{N})
\]

为由交换同构推出的同构

\[
\mathrm{A} \otimes_ {k} \mathrm{A} ^ {\prime} \rightarrow \mathrm{A} ^ {\prime} \otimes_ {k} \mathrm{A}, \quad \mathrm{M} \otimes_ {k} \mathrm{M} ^ {\prime} \rightarrow \mathrm{M} ^ {\prime} \otimes_ {k} \mathrm{M}, \quad \mathrm{N} \otimes_ {k} \mathrm{N} ^ {\prime} \rightarrow \mathrm{N} ^ {\prime} \otimes_ {k} \mathrm{N}, \quad \mathrm{Q} \otimes_ {k} \mathrm{Q} ^ {\prime} \rightarrow \mathrm{Q} ^ {\prime} \otimes_ {k} \mathrm{Q}.
\]

证明这些公式。

\[
\begin{array}{l l} \sigma (a \top a ^ {\prime}) = (- 1) ^ {p q} a ^ {\prime} \top a & \text { for } \quad a \in \operatorname{Tor} _ {p} ^ {A} (Q, M), \quad a ^ {\prime} \in \operatorname{Tor} _ {q} ^ {A ^ {\prime}} (Q ^ {\prime}, M ^ {\prime}) \\ \tau (b \vee b ^ {\prime}) = (- 1) ^ {p q} b ^ {\prime} \vee b & \text { for } \quad b \in \operatorname{Ext} _ {A} ^ {p} (M, N), \quad b ^ {\prime} \in \operatorname{Ext} _ {A ^ {\prime}} ^ {q} (M ^ {\prime}, N ^ {\prime}). \end{array}
\]

\begin{enumerate}
\def\labelenumi{\alph{enumi})}
\setcounter{enumi}{4}
\tightlist
\item
  设 \((\mathcal{E}):0\to \mathbf{M}_2\to \mathbf{M}_1\to \mathbf{M}_0\to 0\) 是 A-模的一个正合列，使得序列
\end{enumerate}

\[
(\mathcal {E} \otimes \mathrm{M} ^ {\prime}): 0 \rightarrow \mathrm{M} _ {2} \otimes_ {k} \mathrm{M} ^ {\prime} \rightarrow \mathrm{M} _ {1} \otimes_ {k} \mathrm{M} ^ {\prime} \rightarrow \mathrm{M} _ {0} \otimes_ {k} \mathrm{M} ^ {\prime} \rightarrow 0
\]

是正合的。设 δ（相应地，Δ）表示连接同态 δ(Q, ℰ)（相应地，δ(Q ⊗ₖ Q′, ℰ ⊗ M′)）。证明对于 a ∈ Tor\^{}A(Q, M₀)，a′ ∈ Tor\^{}A(Q', M')，有 (δa) ⊤ a′ = Δ(a ⊗ a′)。对乘积 ∨ 证明类似的公式。

\begin{enumerate}
\def\labelenumi{\alph{enumi})}
\setcounter{enumi}{5}
\tightlist
\item
  如果 k 是一个域，证明 \(\top\) 是分次 k-模之间的同构；此外，若 A 和 A' 是左诺特的，且 A-模 M 与 A'-模 M' 是有限生成的，则对 ∨ 也有同样的结论。g) 若将模 \(\operatorname{Tor}_{n}^{\mathbf{A}}(\mathbf{Q},\mathbf{M})\) 与习题 6 中定义的集合 \(\operatorname{T}_{n}(\mathbf{Q},\mathbf{M})\) 等同起来，证明公式 \((\mathbf{P},p,q)\top(\mathbf{P}',p',q')=(\mathbf{P}\otimes_{k}\mathbf{P}',p\otimes p',q\otimes q')\) 。
\end{enumerate}

\begin{enumerate}
\def\labelenumi{\arabic{enumi})}
\setcounter{enumi}{7}
\tightlist
\item
  我们沿用习题 7 的记号；我们假设 \(k\) -模 A 与 A' 是投射的。
\end{enumerate}

\begin{enumerate}
\def\labelenumi{\alph{enumi})}
\item
  我们假设有 \(\mathrm{Tor}_{i}^{k}(\mathbf{M},\mathbf{M}^{\prime}) = 0\) （其中 \(i > 0\) ）。证明映射 \(b\mapsto b\vee 1_{\mathbf{M}'}\) 定义了一个零次的分次 \(k\) -同态，从 \(\operatorname{Ext}_{\mathbf{A}}(\mathbf{M},\mathbf{N})\) 到 \(\operatorname{Ext}_{\mathbf{B}}(\mathbf{M}\otimes_{k}\mathbf{M}',\mathbf{N}\otimes_{k}\mathbf{M}')\) ，它在零次诱导出典范同态 \(b\mapsto b\otimes 1_{\mathbf{M}'}\) 。
\item
  我们假设 k-模 M' 是平坦的。设
\end{enumerate}

\[
0 \rightarrow \mathrm{N} \rightarrow \mathrm{R} _ {n} \rightarrow \dots \rightarrow \mathrm{R} _ {1} \rightarrow \mathrm{M} \rightarrow 0
\]

是 A-模的一个正合列，其类为 \(b \in \operatorname{Ext}_{\mathbf{A}}^{n}(\mathbf{M}, \mathbf{N})\) 。证明 B-模的正合列

\[
0 \rightarrow \mathrm{N} \otimes_ {k} \mathrm{M} ^ {\prime} \rightarrow \mathrm{R} _ {n} \otimes_ {k} \mathrm{M} ^ {\prime} \rightarrow \dots \rightarrow \mathrm{R} _ {1} \otimes_ {k} \mathrm{M} ^ {\prime} \rightarrow \mathrm{M} \otimes_ {k} \mathrm{M} ^ {\prime} \rightarrow 0
\]

具有类 \(b \vee 1_{\mathbf{M}'}\) （在 \(\operatorname{Ext}_{\mathbf{B}}^{n}(\mathbf{M} \otimes_{k} \mathbf{M}', \mathbf{N} \otimes_{k} \mathbf{N}')\) 中）。

\begin{enumerate}
\def\labelenumi{\alph{enumi})}
\setcounter{enumi}{2}
\tightlist
\item
  我们假设 k-模 M、M'、N、N' 都是平坦的。证明我们有：
\end{enumerate}

\[
b \vee b ^ {\prime} = (b \vee 1 _ {N ^ {\prime}}) \circ (1 _ {M} \vee b ^ {\prime}) = (- 1) ^ {p q} (1 _ {N} \vee b ^ {\prime}) \circ (b \vee 1 _ {M ^ {\prime}})
\]

对于 \(b \in \operatorname{Ext}_{\mathbf{A}}^{p}(\mathbf{M}, \mathbf{N})\) ， \(b' \in \operatorname{Ext}_{\mathbf{A}'}^{q}(\mathbf{M}', \mathbf{N}')\) 。

\begin{enumerate}
\def\labelenumi{\arabic{enumi})}
\setcounter{enumi}{8}
\tightlist
\item
  我们假设环 A 是交换的。设 B、C 为两个结合的、含单位元的 A-代数；我们用 \(m_{\mathbf{B}}: \mathbf{B} \otimes_{\mathbf{A}} \mathbf{B} \to \mathbf{B}\) 表示 B 的乘法，用 \(m_{\mathbf{C}}\) 表示 C 的乘法。
\end{enumerate}

\begin{enumerate}
\def\labelenumi{\alph{enumi})}
\tightlist
\item
  证明同态
\end{enumerate}

\[
\operatorname{Tor} \left(m _ {\mathrm{B}}, m _ {\mathrm{C}}\right) \circ \top : \operatorname{Tor} ^ {\mathrm{A}} (\mathrm{B}, \mathrm{C}) \otimes_ {\mathrm{A}} \operatorname{Tor} ^ {\mathrm{A}} (\mathrm{B}, \mathrm{C}) \rightarrow \operatorname{Tor} ^ {\mathrm{A}} (\mathrm{B}, \mathrm{C})
\]

（参见习题 7）在 Tor \(^{A}\) (B, C) 上定义一个分次 A-代数的结构，结合且有单位。b) 设 S 是一个微分分次 A-代数，并且 \(p : S_{0} \to B\) 是 A-代数的一个同态，使得 (S, p) 是 B 的一个自由分解（参见 X，第 183 页，习题 18）。证明这个同构

\[
\psi (\mathbf {S}, \mathbf {C}): \operatorname{Tor} ^ {\mathbf {A}} (\mathbf {B}, \mathbf {C}) \rightarrow \mathbf {H} (\mathbf {S} \otimes_ {\mathbf {A}} \mathbf {C})
\]

是 A-代数的同构，如果把 H(S ⊗A C) 赋予由 S ⊗A C 的微分分次代数结构导出的代数结构（X，第 183 页，习题 18，b) 和 c)）。

\begin{enumerate}
\def\labelenumi{\alph{enumi})}
\setcounter{enumi}{2}
\tightlist
\item
  若 A-代数 B 和 C 是交换的，证明分次代数 Tor \(^{A}\) (B, C) 是交错的（参见 X，第 183 页，习题 19）。
\end{enumerate}

\begin{enumerate}
\def\labelenumi{\arabic{enumi})}
\setcounter{enumi}{9}
\tightlist
\item
  设 B 是一个 k-双代数（III，第 148 页），我们经由余单位 \(\beta: \mathbf{B} \to k\) 将其视为一个增广 k-代数（X，第 196 页，习题 27）；我们假设 \(k\) -模 B 是投射的。
\end{enumerate}

\begin{enumerate}
\def\labelenumi{\alph{enumi})}
\tightlist
\item
  设 M、N 为两个左 B-模。习题 7 中的同态 ∨ 等同于一个次数为零的分次同态
\end{enumerate}

\[
\mathrm{H} ^ {\bullet} (\mathbf {B}, \gamma ; \mathbf {M}) \otimes_ {k} \mathrm{H} ^ {\bullet} (\mathbf {B}, \gamma ; \mathbf {N}) \rightarrow \mathrm{H} ^ {\bullet} (\mathbf {B} \otimes_ {k} \mathbf {B}, \gamma \otimes \gamma ; \mathbf {M} \otimes_ {k} \mathbf {N});
\]

通过与由余乘法诱导的同态复合，由此推出一个次数为零的分次同态 \(\cup: H'(B, \gamma; M) \otimes_k H'(B, \gamma; N) \to H'(B, \gamma; M \otimes_k N)\) 。

\begin{enumerate}
\def\labelenumi{\alph{enumi})}
\setcounter{enumi}{1}
\item
  设 C 为结合含幺 k-代数，我们赋予其由 γ 导出的 B-模结构。证明同态 ∪ 在 H'(B, γ ; C) 上定义了结合含幺分次 k-代数的结构；若 C 交换且双代数 B 余交换，则代数 H'(B, γ ; C) 是反交换的。
\item
  我们取 C = k。证明乘积 ∪ 与 \(\operatorname{Ext}_{\mathbf{B}}(k, k)\) 上的复合乘积一致。d) 设 G 是一个群，M 和 N 是两个 \(\mathbf{Z}^{(\mathbf{G})}\) -模。我们赋予 M ⊗\_z N 一个 \(\mathbf{Z}^{(\mathbf{G})}\) -模结构，它由 \(g(x \otimes y) = gx \otimes gy\) （\(g \in G\) ， \(x \in M\) ， \(y \in N\) ）所定义。证明映射
\end{enumerate}

\[
\cup : \mathrm{H} ^ {\bullet} (\mathrm{G}, \mathrm{M}) \otimes_ {\mathbf {z}} \mathrm{H} ^ {\bullet} (\mathrm{G}, \mathrm{N}) \rightarrow \mathrm{H} ^ {\bullet} (\mathrm{G}, \mathrm{M} \otimes_ {\mathbf {z}} \mathrm{N})
\]

是由次数为零的分次同态在同调上诱导的同态

\[
u: \mathrm{C} ^ {*} (\mathrm{G}, \mathrm{M}) \otimes_ {\mathbf {z}} \mathrm{C} ^ {*} (\mathrm{G}, \mathrm{N}) \rightarrow \mathrm{C} ^ {*} (\mathrm{G}, \mathrm{M} \otimes_ {\mathbf {z}} \mathrm{N})
\]

这个同态由 \(u(f \otimes h)(g_1, \ldots, g_{p+q}) = f(g_1, \ldots, g_p) \otimes g_1 \ldots g_p h(g_{p+1}, \ldots, g_{p+q})\) 定义，其中

\[
f \in \mathbf {C} ^ {p} (\mathbf {G}, \mathbf {M}), \quad h \in \mathbf {C} ^ {q} (\mathbf {G}, \mathbf {N}), \quad g _ {1} \dots g _ {p + q} \in \mathbf {G}.
\]

（利用第 185 页习题 5。）

¶ 11) 假设双代数 B 允许一个逆 i（参见 III，第 198 页，习题 4）。设 M 和 N 为 B-模。k-模 M ⊗k N（相应地，Homk(M, N)）具有 B ⊗k B-模（相应地，B ⊗k B°-模）的自然结构；我们用由同态 c（相应地，(1B ⊗ i) ∘ c）导出的 B-模结构赋予它。

\begin{enumerate}
\def\labelenumi{\alph{enumi})}
\tightlist
\item
  证明同态 u（习题 10）定义了一个次数为零的分次同态
\end{enumerate}

\[
\cup : \mathrm{H} ^ {*} (\mathbf {B}, \gamma ; \operatorname{Hom} _ {k} (\mathbf {M}, \mathbf {N})) \otimes_ {k} \mathrm{H} ^ {*} (\mathbf {B}, \gamma ; \mathbf {M}) \rightarrow \mathrm{H} ^ {*} (\mathbf {B}, \gamma ; \mathbf {N}).
\]

\begin{enumerate}
\def\labelenumi{\alph{enumi})}
\setcounter{enumi}{1}
\tightlist
\item
  设 \(0 \to M \to R_n \to \cdots \to R_1 \to k \to 0\) 是 B-模的一个正合列，其中 \(R_i\) 是投射的（\(1 \leqslant i \leqslant n\)）；设 \(\theta\) 为其在 \(\operatorname{Ext}_B^n(k, M) = H^n(B, \gamma; M)\) 中的类。证明对于 \(p \geqslant 1\) ，映射 \(x \mapsto x \cup \theta\) 是一个 \(k\) -同构，从 \(H^p(B, \gamma; \operatorname{Hom}_k(M, N))\) 到 \(H^{p+n}(B, \gamma; N)\) 上。
\end{enumerate}

（证明这个映射经由

\[
\mathrm{H} ^ {p} (\mathrm{B}, \gamma ; \operatorname{Hom} _ {k} (\mathrm{M}, \mathrm{N})) \rightarrow \operatorname{Ext} _ {\mathrm{B}} ^ {p} (\mathrm{M}, \mathrm{N}) \xrightarrow {\theta} \operatorname{Ext} _ {\mathrm{B}} ^ {p + n} (k, \mathrm{N}).
\]

\begin{enumerate}
\def\labelenumi{\alph{enumi})}
\setcounter{enumi}{2}
\tightlist
\item
  设 G 为一个群，N 为一个 \(\mathbf{Z}^{(G)}\) -模。设 F 为一个自由群， \(p: F \to G\) 一个满同态，R = Ker( \(p\) ), \(\overline{R} = R/(R, R)\) , \(\overline{F} = F/(R, R)\) 。该扩张 \(\overline{R} \to \overline{F} \to G\) 定义了一个类 \(\varepsilon \in H^2(G, \overline{R})\) ；证明映射 \(x \mapsto x \cup \varepsilon\) 从 \(H^p(G, \operatorname{Hom}_{\mathbf{Z}}(\overline{R}, N))\) 到 \(H^{p+2}(G, N)\) 是一个同构（\(p \geqslant 1\) ；利用习题 3，d)，第 179 页）。
\end{enumerate}

